\section{A recursion category from a three-point space}\label{sec:cosheaves}

We construct a recursion category $\mathscr R_X$ from set-valued
cosheaves in the sense of Curry~\cite[Definition~2.6]{CurryCosheaves}
on a three-point space $X$. A cosheaf on $X$ is determined up to isomorphism
by a diagram of shape $J=(1\leftarrow 0\to 2)$. Evaluating an internal propositional theory
of $\operatorname{Kar}(\mathscr R_X)$ at either endpoint gives a consistent theory.
The indices of its provable atoms and those of its refutable atoms
contain, respectively, the two members of a recursively inseparable pair.
In that endpoint language, every consistent extension with a computably
enumerable axiom set leaves an atom undecided. The internal first-order theory constructed
in Section~\ref{subsec:fork-first-order} restricts to this propositional
theory.

For a topological space $X$, let $\mathcal O(X)$ be the category
of open sets and inclusions. A \emph{set-valued cosheaf} is a
covariant functor $C:\mathcal O(X)\to\mathbf{Set}$ for which
every open cover $(O_\lambda)_\lambda$ of $O$ gives a coequalizer
\[
 \coprod_{\lambda,\mu} C(O_\lambda\cap O_\mu)
 \rightrightarrows \coprod_\lambda C(O_\lambda)
 \longrightarrow C(O).
\]
The two parallel morphisms are induced by the inclusions into the two
members of the cover; the empty cover requires
$C(\varnothing)=\varnothing$.
The \emph{Sierpi\'nski space} has two points and exactly three open
sets: the empty set, one singleton, and the whole space.

Let $X=\{a,b,c\}$ have opens
$\varnothing,\{c\},\{a,c\},\{b,c\},X$.
The open subspaces $\{a,c\}$ and $\{b,c\}$ are Sierpi\'nski spaces
with common open singleton $\{c\}$. Thus $X$ is obtained by gluing
two copies of the Sierpi\'nski space along that singleton.
Its category of set-valued cosheaves is equivalent to
$\mathbf{Set}^{J}$, where $J$ has objects $0,1,2$ and nonidentity
morphisms $0\to1$ and $0\to2$.
The equivalence sends a cosheaf $C$ to the functor with values
$C(\{c\}),C(\{a,c\}),C(\{b,c\})$ at $0,1,2$, respectively,
and morphisms induced by the inclusions of these open sets.
The cosheaf value $C(X)$ is the pushout of
$C(\{a,c\})\leftarrow C(\{c\})\to C(\{b,c\})$
(Proposition~\ref{prop:fork-cosheaves}).

Use the category $\mathcal C_{\mathrm{rec}}$ of
Section~\ref{subsec:numerical-realization}. Its objects are pairs
$C=(m,C_0)$ with carrier $|C|=C_0$, and its morphisms are total
computable functions on the carriers. For a c.e.\ subset
$C_0\subseteq\N^m$, writing $C_0$ for the object $(m,C_0)$
abbreviates this pair. Membership, subset and image expressions,
and elementwise computations involving numerical objects refer
to their carriers.

Put $\mathcal E=\mathcal C_{\mathrm{rec}}^J$.
An object is a covariant functor $Y:J\to\mathcal C_{\mathrm{rec}}$,
given by the objects $Y_j=Y(j)$, for $j=0,1,2$, and the
morphisms $r_i^Y=Y(0\to i):Y_0\to Y_i$, for $i=1,2$.
A morphism $f:Y\to Z$ is a natural transformation: its components
are morphisms $f_j:Y_j\to Z_j$ of $\mathcal C_{\mathrm{rec}}$ satisfying
$r_i^Zf_0=f_i r_i^Y$.
A \emph{c.e.\ subfunctor} $D\subseteq Y$ has component objects
$D_j=(m_j,D_{j,0})$, where $Y_j=(m_j,Y_{j,0})$ and
$D_{j,0}\subseteq|Y_j|$ is c.e.\ for each $j\in\{0,1,2\}$.
The conditions $|r_i^Y|(|D_0|)\subseteq|D_i|$, for $i=1,2$,
define its morphisms $r_i^D:D_0\to D_i$ by restriction.
The component inclusions give a natural transformation $D\hookrightarrow Y$.
The underlying sets and functions define a functor
$\mathcal E\to\mathbf{Set}^J$. Composing with the quasi-inverse of
the cosheaf equivalence above assigns a set-valued cosheaf to each
object of $\mathcal E$ and a natural transformation to each morphism.

Put
$\mathcal A=\operatorname{Par}(\mathcal E,\operatorname{Mono})$.
Monomorphisms of $\mathcal E$ are componentwise injections.
Thus a partial morphism is equivalently a computably enumerable
subfunctor $D\subseteq Y$ together with a natural transformation
$D\to Z$. Its restriction is the partial identity on $D$, and its range is the
partial identity on the subfunctor of component images.
Proposition~\ref{prop:fork-structure} proves these assertions.

Define an object $V:J\to\mathcal C_{\mathrm{rec}}$ of $\mathcal E$
by the following objects $V_j=V(j)$ and morphisms
$r_i^V=V(0\to i)$:
\begin{equation}\label{eq:universal-fork}
 \begin{gathered}
 V_0=\N^4,\qquad V_1=V_2=\N^2,\\
 r_i^V:V_0\to V_i,\qquad r_i^V(k,z,x_1,x_2)=(k,x_i)
 \quad(i=1,2).
 \end{gathered}
\end{equation}
For $A=(m,A_0)$ in $\mathcal C_{\mathrm{rec}}$, put
$m(A)=m$. Thus $|A|\subseteq\N^{m(A)}$,
$m(V_0)=4$, and $m(V_1)=m(V_2)=2$.
Fix recursive numerical encodings of nonempty finite sequences of tuples,
\[
 \operatorname{code}_m:
 \coprod_{n\ge1}(\N^m)^n\xrightarrow{\cong}W_m\subseteq\N
 \qquad(m\ge0),
\]
where each $W_m$ is recursive and the length and entry operations
are recursive. For an object $A$ of $\mathcal C_{\mathrm{rec}}$, put
\[
 A^\#=\{\operatorname{code}_{m(A)}(a_1,\ldots,a_n):
                 n\ge1,\ a_1,\ldots,a_n\in A\}\subseteq W_{m(A)}.
\]
This is a c.e.\ numerical carrier. For each $Y\in\mathcal E$, let
$(Y^\#)_j=(Y_j)^\#$. Its morphism
$(Y^\#)(0\to i):(Y_0)^\#\to(Y_i)^\#$ applies $r_i^Y$ to each
entry of a word, for $i=1,2$.
Choose isomorphisms in $\mathcal E$
$V\cong V\times V$, $V\cong V+V$, and $V\cong V^\#$.
Together with their inverses, these form a \emph{frame} at $V$
\cite[Section 3]{Heller}.
Let $\mathscr R_X$ denote the full subcategory of $\mathcal A$
on the objects obtained from $V$ by finitely many applications of
the chosen binary restriction products, binary coproducts, and
nonempty-word formation.
The isomorphisms, constructed in Lemma~\ref{lem:fork-frames}, make
$\mathscr R_X$ isotypical.

Theorem~\ref{thm:fork-turing} constructs a Turing morphism
$\tau:V\times V\rightharpoonup V$ in $\mathcal A$.
For every object $Y$ of $\mathcal A$ and every morphism
$F:Y\times V\rightharpoonup V$ in $\mathcal A$, there is a
total morphism $p_F:Y\to V$ such that $\tau(p_F\times1_V)=F$.
By fullness this property holds in $\mathscr R_X$ when
$Y$ belongs to $\mathscr R_X$. Together with the inherited
prodominical structure, it makes $\mathscr R_X$ a recursion category.

Every object $Y$ of $\mathcal E$ admits a total morphism $s_Y:Y\to V$ in
$\mathcal A$ with a partial inverse $r:V\rightharpoonup Y$
such that $r s_Y=1_Y$ and $s_Y r=\bar r$
(Lemma~\ref{lem:fork-dominance}). Thus $V$ is a
\emph{universal object} of $\mathcal A$.
These embeddings give the equivalence
$\operatorname{Split}(\mathscr R_X)\simeq\mathcal A$
proved in Proposition~\ref{prop:fork-recursion}.
Every idempotent in $\mathcal A$ splits through its image.
Corollary~\ref{cor:fork-karoubi} gives a restriction equivalence
$Q:\operatorname{Kar}(\mathscr R_X)\to\mathcal A$. We construct the
formula and derivation objects in $\mathcal E$, canonically isomorphic to
$\operatorname{Tot}(\mathcal A)$.
A quasi-inverse to $Q$ transports them to the total category of
$\operatorname{Kar}(\mathscr R_X)$. At an endpoint $i$, the evaluation
of a transported object is given, up to the counit isomorphism, by
its original $i$th component.

\subsection{Complementary restriction idempotents and inseparability}

Let $e,e'$ be complementary restriction idempotents on $V$ in
$\mathscr R_X$. Write $\iota_C:C\hookrightarrow V$ and
$\iota_{C'}:C'\hookrightarrow V$ for the inclusions in $\mathcal E$
through which $e,e'$ split in $\mathcal A$.
The component carriers $C_j,C'_j$ are disjoint c.e.\ subsets whose union
is $V_j$ for each $j\in\{0,1,2\}$. Consequently the copairing
$\chi=[\iota_C,\iota_{C'}]:C+C'\to V$ has bijective components
$\chi_j:C_j+C'_j\to V_j$. Each inverse is
computable by running the membership searches for $C_j$ and $C'_j$ in
parallel and returning the element in the corresponding coproduct
summand. Naturality of $\chi$ gives
\[
 r_i^{C+C'}\chi_0^{-1}=\chi_i^{-1}r_i^V:V_0\longrightarrow C_i+C'_i
 \qquad(i=1,2).
\]
Thus these inverse components define $\chi^{-1}:V\to C+C'$ in
$\mathcal E$, proving that $\chi$ is an isomorphism.
For the injections $i_0,i_1:V\to V+V$, define
$\ell_e=[i_1\iota_C,i_0\iota_{C'}]\chi^{-1}:V\to V+V$.
It satisfies $\ell_e e=i_1e$ and $\ell_e e'=i_0e'$ and belongs to
$\mathscr R_X$ by fullness. The full inclusion
$I:\mathscr R_X\hookrightarrow\mathcal A$ is a range functor,
$I(1_V)\ne0$, and inverse images along total morphisms in $\mathcal A$ preserve
binary joins by Proposition~\ref{prop:fork-structure}.
Corollary~\ref{cor:source-conditionals} therefore gives an inseparable
pair in $\mathscr R_X$.
The category $\mathscr R_X$ is not restriction-equivalent to the
recursion category of partial recursive functions
(Proposition~\ref{prop:fork-distinction}).

\subsection{Evaluation functors}

Let
$\mathcal K=\operatorname{Par}(\mathcal C_{\mathrm{rec}},\operatorname{Mono})$.
Its objects are c.e.\ numerical carriers, and its morphisms are partial
recursive functions with c.e.\ domains contained in their sources.
The functor
$\mathsf Q_{\mathrm{num}}:\operatorname{Split}(\mathrm{Rec})\to\mathcal K$
constructed in Proposition~\ref{prop:numerical-structure} is a
restriction equivalence.
For $j\in\{0,1,2\}$, component evaluation defines
\begin{equation}\label{eq:endpoint-functor}
 E_j:\mathcal A\longrightarrow\mathcal K,\qquad
 E_j(Y)=Y_j,\qquad
 E_j(Y\xleftarrow{\iota_D}D\xrightarrow{f}Z)
   =(Y_j\xleftarrow{(\iota_D)_j}D_j\xrightarrow{f_j}Z_j).
\end{equation}
For the Turing morphism $\tau:V\times V\rightharpoonup V$ of
Theorem~\ref{thm:fork-turing}, write
$\tau_j=E_j(\tau):V_j\times V_j\rightharpoonup V_j$.

\begin{proposition}\label{prop:endpoint}
The covariant functors $E_j$ preserve restrictions, zeros, ranges,
binary restriction products and coproducts, and finite joins of
restriction idempotents.
For $i=1,2$, the Turing morphism
$\tau:V\times V\rightharpoonup V$ constructed in
Theorem~\ref{thm:fork-turing} is sent
to a Turing morphism on $\N^2$ in $\mathcal K$.
\end{proposition}
\begin{proof}
Component evaluation preserves identities, monomorphisms, and
pullbacks in the total categories. Consequently \eqref{eq:endpoint-functor}
is independent of the representative span and preserves partial-map
composition. The stated operations are represented componentwise by
partial identities, empty sets, images, products, disjoint unions, and
finite unions, respectively.

Use the recursive bijection $\operatorname{encode}_2:\N^2\to\N$,
its inverse $\operatorname{decode}_2:\N\to\N^2$, and the acceptable numbering
$(\varphi_p)_{p\in\N}$ chosen in the construction of $\tau$ in
Appendix~\ref{app:fork-structure}. For $i=1,2$, the component $\tau_i$ is given by
\eqref{eq:fork-turing-endpoint}.
For a partial recursive map
$F:Z\times\N^2\rightharpoonup\N^2$ with c.e.\ carrier $Z$,
the parameter theorem gives a total recursive map $s:Z\to\N$ such that
\[
 \varphi_{s(z)}
 =\operatorname{encode}_2\circ F(z,-)\circ\operatorname{decode}_2
 :\N\rightharpoonup\N
 \qquad(z\in Z).
\]
Thus $z\mapsto(0,s(z)):Z\to\N^2$ is the required total parameter morphism.
\end{proof}

For $i\in\{1,2\}$, let $H_i:J\to\mathcal C_{\mathrm{rec}}$
have component $\{*\}$ at $i$ and empty components elsewhere.
For $Y\in\mathcal E$ and $y\in Y_i$, let $p_y:H_i\to Y$ be the total morphism
whose component at $i$ sends $*$ to $y$. Every total morphism
$H_i\to Y$ has this form for a unique $y\in Y_i$. A partial
morphism from $H_i$ is either total or zero, since $H_i$ has
two restriction idempotents. Postcomposition therefore defines
$\Phi_i:\mathcal A\to\Pfn$ by
\[
 \Phi_i(Y)=\operatorname{Tot}(\mathcal A)(H_i,Y),\qquad
 \Phi_i(f)(p)=fp\quad\hbox{when }fp\hbox{ is total}.
\]
If $fp$ is not total, it is zero, and every further composite
remains zero. If it is total, further postcomposition has exactly
the domain specified by the next partial function. This proves
$\Phi_i(gf)=\Phi_i(g)\Phi_i(f)$, and $\Phi_i(1_Y)(p)=p$.
Let $\lvert-\rvert:\mathcal K\to\Pfn$ send a c.e.\ numerical carrier to its
underlying set and a partial recursive morphism to the same partial
function. The bijections $p_y\mapsto y$ give a natural isomorphism
$\Phi_i\cong\lvert-\rvert\circ E_i$.
Indeed, $fp_y=p_{E_i(f)(y)}$ when $E_i(f)(y)$ is defined, and $fp_y=0$
otherwise.
For a restriction idempotent $e$ on $Y$, put
\begin{equation}\label{eq:probe-domain}
 |e|_i:=\{y\in Y_i:ep_y=p_y\}
       =\{y\in Y_i:E_i(e)(y)=y\}.
\end{equation}
This is the c.e.\ subset on which $E_i(e)$ is the partial identity.

The morphism $E_0(\tau)$ is not a Turing morphism. For the parameter
object $1$, its universal property would represent each partial recursive
$F:V_0\rightharpoonup V_0$ by a total point $1\to V_0$, using
$1\times V_0\cong V_0$. The total function $F:V_0\to V_0$ given by
$F(k,z,x_1,x_2)=(0,0,z,0)$ cannot equal $\tau_0(p,-)$ for any $p\in V_0$.
Let $D_\tau\hookrightarrow V\times V$ be the domain subfunctor of
$\tau$ constructed in Appendix~\ref{app:fork-turing}.
If $F=\tau_0(p,-)$, totality of $F$ gives $(p,x)\in D_{\tau,0}$
for every $x\in V_0$. Naturality of $D_\tau\to V$ then gives
\[
 r_1^V F(x)=\tau_1(r_1^Vp,r_1^Vx).
\]
The right-hand side is unchanged when only the $z$ coordinate
of $x=(k,z,x_1,x_2)$ changes, whereas the left-hand side is $(0,z)$.

\subsection{Propositional theories and component theorem sets}

Let $\zeta:V\to V+V$ be the coproduct isomorphism of
Lemma~\ref{lem:fork-frames}. At each component $j$, the map
$\zeta_j$ selects the summand according to the parity of the
first coordinate in $V_j$. Define complementary restriction idempotents
$\varepsilon^V_0,\varepsilon^V_1:V\rightharpoonup V$ by
\[
 \varepsilon^V_0=\zeta^{-1}(1_V+0_V)\zeta,\qquad
 \varepsilon^V_1=\zeta^{-1}(0_V+1_V)\zeta.
\]
Put $d=\tau\Delta_V:V\rightharpoonup V$. Let $A,B\subseteq V$ be the
c.e.\ subfunctors on which the restriction idempotents
$d^*\varepsilon^V_0,d^*\varepsilon^V_1$ are the partial identities
$1_A,1_B$, respectively. These idempotents are inseparable by
Corollary~\ref{cor:source-conditionals} for the full inclusion
$\mathscr R_X\hookrightarrow\mathcal A$. They remain inseparable
in $\mathcal A$: fullness leaves the restriction idempotents on $V$, their order,
and their complements unchanged.
Write $A_j=A(j)$ and $B_j=B(j)$
for $j=0,1,2$. For $i\in\{1,2\}$, put $d_i=E_i(d)$.
By Proposition~\ref{prop:endpoint},
\begin{align*}
 A_i&=\{x:d_i(x)\text{ is defined with even first coordinate}\},\\
 B_i&=\{x:d_i(x)\text{ is defined with odd first coordinate}\}.
\end{align*}
Let $\mathcal K_2$ be the full subcategory of $\mathcal K$ on the
objects obtained from $\N^2$ by the chosen binary restriction
products and coproducts. Every object of $\mathcal K_2$ is recursively
isomorphic to $\N^2$. The restriction products, coproducts, and
prodominical zero law restrict to this subcategory. By fullness,
$E_i\tau$ and the total parameter morphisms of
Proposition~\ref{prop:endpoint} belong to $\mathcal K_2$ whenever
their sources do. Thus $\mathcal K_2$ is a recursion category.

On $\mathcal K_2$, the underlying-set functor preserves
restrictions, zeros, and ranges, since it sends each morphism to its
underlying partial function. It also preserves the chosen products
and coproducts, so it is a range functor.
For complementary restriction idempotents $e,e':W\rightharpoonup W$
in $\mathcal K_2$, the construction in
Section~\ref{sec:numerical-incompleteness} gives a total recursive
morphism $\ell_e:W\to W+W$ with
$\ell_e e=i_1e$ and $\ell_e e'=i_0e'$.
It belongs to $\mathcal K_2$ by fullness. With the componentwise coproduct
identification, $E_i\zeta:\N^2\to\N^2+\N^2$ is a coproduct
isomorphism in $\mathcal K_2$.
Since the identity on $\N^2$ is nonzero and inverse images along
total morphisms in $\Pfn$ preserve unions,
Corollary~\ref{cor:source-conditionals} makes the partial identities
on $A_i,B_i$ inseparable in $\mathcal K_2$. A recursive separator
would define a complemented restriction idempotent on $\N^2$ in
this full subcategory. Hence $A_i,B_i$ are recursively inseparable.

The category $\mathcal E$ is regular with finite coproducts, and its
images are componentwise images (Proposition~\ref{prop:fork-total}).
Apply the propositional construction of
Appendix~\ref{app:propositional-theories} in this category, with atom
object $V$ and additional axioms determined by $A,B\subseteq V$.

For $j=0,1,2$, let $S_j$ be the set of finite rooted ordered formula trees
with constructors
\[
 \theta_j(x)\quad(x\in V_j),\qquad
 \nu_j(\varphi)\quad(\varphi\in S_j),\qquad
 \varphi\Rightarrow_j\psi\quad(\varphi,\psi\in S_j).
\]
The constructor images are pairwise disjoint. Recursive encodings with
recursive parsing make $S_j$ a recursive numerical carrier, since
$V_j$ is recursive. For $i=1,2$, replacing every atom label $x$ by
$r_i^V(x)$ and retaining the constructors gives a total recursive map
$\sigma_i:S_0\to S_i$. These maps define a functor
$S:J\to\mathcal C_{\mathrm{rec}}$. The constructors commute with
relabeling, so they define morphisms $\theta:V\to S$,
$\nu:S\to S$, and ${\Rightarrow}:S\times S\to S$ of $\mathcal E$.
Define $P:\mathcal E\to\mathcal E$ by
$P(Z)=V+Z+(Z\times Z)$ and
$P(f)=1_V+f+(f\times f)$. Its constructor algebra is
$\iota=[\theta,\nu,\Rightarrow]:P(S)\to S$.

\begin{proposition}\label{prop:fork-syntax-initial}
The algebra $\iota$ is initial: for every
$\alpha:P(Z)\to Z$ in $\mathcal E$, there is a unique morphism
$u:S\to Z$ in $\mathcal E$ satisfying $u\iota=\alpha P(u)$.
\end{proposition}
\begin{proof}
At component $j$, evaluate a formula tree from its leaves to its root,
using the operations given by $\alpha_j:V_j+Z_j+Z_j^2\to Z_j$.
Induction proves termination, membership of the result in $Z_j$, and
the algebra equation. Evaluation is partial recursive and total on
$S_j$, so it defines a morphism $u_j:S_j\to Z_j$ of
$\mathcal C_{\mathrm{rec}}$.

Write $z_i=Z(0\to i):Z_0\to Z_i$. Naturality of $\alpha$
says that $z_i$ commutes with each constructor operation, with
$r_i^V$ applied to atom labels. Induction therefore gives
$z_i u_0=u_i\sigma_i$. Thus the $u_j$ form a morphism of
$\mathcal E$. Any algebra morphism agrees with these evaluations
on atoms and then on negations and implications by induction,
proving uniqueness.
\end{proof}

The logical axiom schemes define morphisms
$\Lambda_1,\Lambda_3:S^2\to S$ and $\Lambda_2:S^3\to S$.
For $\alpha,\beta,\gamma\in S_j$, their components are
\begin{align*}
 \Lambda_{1,j}(\alpha,\beta)
   &=\alpha\Rightarrow_j(\beta\Rightarrow_j\alpha),\\
 \Lambda_{2,j}(\alpha,\beta,\gamma)
   &=(\alpha\Rightarrow_j(\beta\Rightarrow_j\gamma))
     \Rightarrow_j((\alpha\Rightarrow_j\beta)
                         \Rightarrow_j(\alpha\Rightarrow_j\gamma)),\\
 \Lambda_{3,j}(\alpha,\beta)
   &=(\nu_j\beta\Rightarrow_j\nu_j\alpha)
                         \Rightarrow_j(\alpha\Rightarrow_j\beta).
\end{align*}
The morphism
$[\theta|_A,\nu\theta|_B]:A+B\to S$ inserting the positive and negative
atomic axioms is monic: each component is
injective, because the constructors are
injective and the atom and negation constructors have disjoint images.
Its image is the subobject of signed axioms. Put
\[
 \begin{gathered}
 \mathrm{Ax}=S^2+S^3+S^2+A+B,\\
 \lambda=[\Lambda_1,\Lambda_2,\Lambda_3,\theta|_A,\nu\theta|_B]:
                 \mathrm{Ax}\longrightarrow S.
 \end{gathered}
\]
Each element of $\mathrm{Ax}_j$ specifies an axiom instance.
The morphism $\lambda_j$ gives its conclusion.

Write $\pi_1,\pi_2:S^2\to S$ for the projections and put
$\mu=\langle\pi_1,\Rightarrow\rangle:S^2\to S^2$.
For $h:Y\to S$ in $\mathcal E/S$, let $M(Y,h)$ be the pullback
\[
\begin{tikzcd}[column sep=large]
 M(Y,h)\arrow[r,"v_h"]\arrow[d,"w_h"'] &
 Y\times Y\arrow[d,"h\times h"]\\
 S\times S\arrow[r,"\mu"'] & S\times S.
\end{tikzcd}
\]
An element $(p_1,p_2,\alpha,\beta)$ of $M(Y,h)_j$ satisfies
$h_j(p_1)=\alpha$ and $h_j(p_2)=\alpha\Rightarrow_j\beta$.
The conclusion morphism is
$\kappa_h=\pi_2w_h:M(Y,h)\to S$. Define
\[
 \mathsf H:\mathcal E/S\longrightarrow\mathcal E/S,
 \qquad
 \mathsf H(Y,h)=
       (\mathrm{Ax}+M(Y,h),[\lambda,\kappa_h]).
\]
For a slice morphism $g:(Y,h)\to(Y',h')$, the pullback gives
$M(g):M(Y,h)\to M(Y',h')$ with
$v_{h'}M(g)=(g\times g)v_h$ and $w_{h'}M(g)=w_h$.
Put $\mathsf H(g)=1_{\mathrm{Ax}}+M(g)$; pullback uniqueness
proves the functor laws. The rule data of
Appendix~\ref{app:closed-subobjects} are
$R_0=\mathrm{Ax}$, $q_0=\lambda$, $R_1=0$, and
$R_2=S^2$, $p_2=\mu$, $q_2=\pi_2$; the morphism $p_0:R_0\to1$
is unique. Since $\mathcal E$ is extensive, the unary pullback
summand is initial, giving the displayed form of $\mathsf H$.

At component $j$, a derivation is a finite rooted ordered tree.
Each leaf is labeled by a member of $\mathrm{Ax}_j$ and concludes
its image under $\lambda_j$. Each nonleaf represents a modus-ponens
inference and has two children concluding $\alpha$ and
$\alpha\Rightarrow_j\beta$; it concludes $\beta$. Let
$\mathrm{Prf}_j$ be the carrier of these trees. Tree shape, logical
parameters, and premise matching are recursively decidable; the
finitely many positive and negative axiom labels require semidecision
of membership in $A_j$ and $B_j$. Dovetailing these tests shows that
$\mathrm{Prf}_j$ is c.e. Reading the conclusion gives a morphism
$c_{\mathrm p,j}:\mathrm{Prf}_j\to S_j$ of $\mathcal C_{\mathrm{rec}}$.

For $i=1,2$, relabel formula parameters by $\sigma_i$ and atomic
axiom labels by $r_i^V$, retaining the rule at each node. Since $A,B$ are
subfunctors, $r_i^V(A_0)\subseteq A_i$ and
$r_i^V(B_0)\subseteq B_i$. Relabeling therefore preserves axiom
leaves. Commutation of $\sigma_i$ with implication preserves premise
matching at modus-ponens nodes. The resulting morphisms
$\rho_i:\mathrm{Prf}_0\to\mathrm{Prf}_i$ satisfy
$c_{\mathrm p,i}\rho_i=\sigma_i c_{\mathrm p,0}$, by induction. They define an object
$\mathrm{Prf}$ of $\mathcal E$, and the $c_{\mathrm p,j}$ define
$c_{\mathrm p}:\mathrm{Prf}\to S$ in $\mathcal E$.

Forming a leaf and adjoining a modus-ponens root give morphisms
$\mathsf{leaf}:\mathrm{Ax}\to\mathrm{Prf}$ and
$\mathsf{mp}:M(\mathrm{Prf},c_{\mathrm p})\to\mathrm{Prf}$ in $\mathcal E$.
They are recursive on their carriers and commute with relabeling.
The injective implication constructor makes the consequent at a
modus-ponens node unique. The equations
$c_{\mathrm p}\,\mathsf{leaf}=\lambda$ and $c_{\mathrm p}\,\mathsf{mp}=\kappa_{c_{\mathrm p}}$ make
$\delta=[\mathsf{leaf},\mathsf{mp}]:\mathsf H(\mathrm{Prf},c_{\mathrm p})\to(\mathrm{Prf},c_{\mathrm p})$
an algebra in $\mathcal E/S$.

\begin{proposition}\label{prop:fork-proof-initial}
The algebra $\delta$ is initial. For every
$\eta:\mathsf H(Y,h)\to(Y,h)$ in $\mathcal E/S$, there is a
unique morphism $u:\mathrm{Prf}\to Y$ of $\mathcal E$ with
$hu=c_{\mathrm p}$ and $u\delta=\eta\mathsf H(u)$.
\end{proposition}
\begin{proof}
Evaluate a derivation from its leaves to its root using $\eta_j$.
At a leaf labeled by $\xi\in\mathrm{Ax}_j$, the value has
$h_j$-image $\lambda_j(\xi)$ because $\eta$ is a slice morphism.
If the children of a modus-ponens node conclude $\alpha$ and
$\alpha\Rightarrow_j\beta$, their evaluated values, together
with $\alpha,\beta$, belong to $M(Y,h)_j$. The corresponding
operation of $\eta_j$ returns a value with $h_j$-image $\beta$.
Induction proves termination and gives a partial recursive function
$u_j:\mathrm{Prf}_j\to Y_j$ total on its source carrier, with
$h_ju_j=c_{\mathrm p,j}$.

Naturality of $\eta$, $\lambda$, and the pullback projections
makes each evaluation step commute with relabeling. Induction gives
$Y(0\to i)u_0=u_i\rho_i$, so the $u_j$ form a morphism
$u:\mathrm{Prf}\to Y$ of $\mathcal E$. Its defining recursion
is the algebra equation. Any other algebra morphism agrees on
leaves and then on each modus-ponens node by induction, proving
uniqueness.
\end{proof}

For a restriction idempotent $t$ on $S$, let $D_t\subseteq S$
be the c.e. subfunctor on which $t$ is the partial identity. View
$c_{\mathrm p}$ as a total morphism of $\mathcal A$ and put $t_0=\ran(c_{\mathrm p})$.

\begin{theorem}\label{thm:fork-proof-presentation}
The subfunctor $D_{t_0}\subseteq S$ is the least c.e. subfunctor
containing the logical and positive and negative atomic axioms and
closed under modus ponens. Its component theories are consistent,
and
\begin{equation}\label{eq:atomic-inequalities}
 1_A\le\theta^*t_0,\qquad
 1_B\le\theta^*\nu^*t_0,\qquad t_0\nu^*t_0=0_S.
\end{equation}
\end{theorem}
\begin{proof}
Apply Theorem~\ref{thm:least-theorem-subobject} to the regular
category $\mathcal E$ and the initial algebra of
Proposition~\ref{prop:fork-proof-initial}. It identifies the image
of $c_{\mathrm p}$ with the least subobject closed under these rules. Images
in $\mathcal E$ are componentwise, so its components are
$(D_{t_0})_j=c_{\mathrm p,j}(\mathrm{Prf}_j)\subseteq S_j$.
Lemma~\ref{lem:range-regular-image}, applied in
$\mathcal A=\Par(\mathcal E,\Mono)$, identifies its partial identity
with $t_0=\ran(c_{\mathrm p})$, and hence the image subfunctor with $D_{t_0}$.
The one-node positive and negative proofs give morphisms
$A\to\mathrm{Prf}$ and $B\to\mathrm{Prf}$ whose composites
with $c_{\mathrm p}$ are $\theta|_A$ and $\nu\theta|_B$. These factorizations
give the first two inequalities in \eqref{eq:atomic-inequalities}.

For consistency, fix $j$ and suppose derivations $p,p'$ conclude
$\varphi$ and $\nu_j\varphi$. Let $L\subseteq A_j$ be the finite
set of labels at positive leaves of either derivation. Assign truth
value $1$ precisely to atoms with labels in $L$, and value $0$ to
all other atoms. Extend this assignment by Boolean negation and
implication. Every positive leaf is true. Every negative leaf has
its label in $B_j$, disjoint from $A_j$ and hence from $L$, so it
is true. The three logical schemes are Boolean tautologies, and
modus ponens preserves truth. Induction on both derivations would
make $\varphi$ and its negation true, a contradiction. Thus
$(D_{t_0})_j\cap\nu_j^{-1}((D_{t_0})_j)=\varnothing$ at every component, which
is $t_0\nu^*t_0=0_S$ because restriction operations in
$\mathcal A$ are componentwise.
\end{proof}

Write $\sigma_i=S(0\to i):S_0\to S_i$ for $i=1,2$.
A compatible family of c.e.\ theory extensions consists of c.e.\ sets
$T_j\subseteq S_j$, for $j=0,1,2$, containing $(D_{t_0})_j$ and closed
under modus ponens, with $\sigma_i(T_0)\subseteq T_i$ for $i=1,2$.
These sets form a subfunctor $T\subseteq S$.
The family is consistent if each theory is consistent, that is,
$T_j\cap\nu_j^{-1}(T_j)=\varnothing$ for every $j$.
Each compatible family of consistent c.e.\ deductively closed extensions
determines a restriction idempotent $t_T:S\rightharpoonup S$, the partial
identity on its theorem subfunctor $T\subseteq S$.
It satisfies $t_0\le t_T$ and $t_T\nu^*t_T=0_S$.
All three theories are complete precisely when
$t_T\vee\nu^*t_T=1_S$, since the $j$th component of $\nu^*t_T$
is the partial identity on $\nu_j^{-1}(T_j)$.

\begin{theorem}\label{thm:cosheaf-observation}
Fix $i\in\{1,2\}$. Let $T\subseteq S$ be a compatible family of
consistent c.e.\ theory extensions as above, and let
$t_T:S\rightharpoonup S$ be its partial identity. The total morphism $E_i\theta$ induces the function
$\theta_i:|V_i|=\N^2\to |S_i|$ on carriers. Define
\[
 P_{i,T}=\{x\in\N^2:\theta_i(x)\in T_i\},\qquad
 N_{i,T}=\{x\in\N^2:\nu_i(\theta_i(x))\in T_i\}.
\]
The restriction idempotents $t_T,\nu^*t_T$ are inseparable in $\mathcal A$.
The sets $P_{i,T},N_{i,T}$ are disjoint c.e.\ subsets of $\N^2$ satisfying
$A_i\subseteq P_{i,T}$ and $B_i\subseteq N_{i,T}$.
They are recursively inseparable. The assignment
$T\mapsto(\omega\times\omega,P_{i,T},N_{i,T})$,
with the identity of $\omega\times\omega$ for extension inclusions, is a
covariant functor from the poset of these families, ordered by componentwise
inclusion, to $\operatorname{Pair}(\mathrm{Rec})$.
The theory with theorem set $T_i$ is incomplete.
\end{theorem}
\begin{proof}
The category $\mathcal A$ has restriction zeros and joins of restriction idempotents
preserved by inverse images along total morphisms. Its restriction idempotents $1_A,1_B$ are inseparable,
and \eqref{eq:atomic-inequalities} gives the hypotheses of
Corollary~\ref{cor:domain-incompleteness}. Hence $t_T,\nu^*t_T$ are
inseparable.

The atom morphism $E_i\theta:\N^2\to S_i$ is total recursive.
The component theorem and refutation sets $|t_T|_i$ and $|\nu^*t_T|_i$
are c.e.; their inverse images under the total recursive map
$E_i\theta$ are therefore c.e. Consistency and preservation of meets and
zero make them disjoint. Applying $E_i$ to
\eqref{eq:atomic-inequalities} and using $t_0\le t_T$ gives
the two inclusions. A recursive separator would consequently
separate $A_i,B_i$, which is impossible.
If $P_{i,T}\cup N_{i,T}=\N^2$, simultaneous enumeration of these
disjoint c.e.\ sets would decide membership in $P_{i,T}$ and hence
separate the pair. Take $x\in\N^2\setminus(P_{i,T}\cup N_{i,T})$.
By the definitions of the two inverse images,
\[
  \theta_i(x)\notin T_i,\qquad \nu_i\theta_i(x)\notin T_i .
\]
Thus an atomic sentence is undecided at endpoint $i$.

If $T\subseteq T'$, monotonicity gives inclusions of both component
sets, so the identity of $\omega\times\omega$ is the required pair morphism.
The functor laws follow from the identity maps.
\end{proof}

\begin{corollary}\label{cor:cosheaf-essential-incompleteness}
For each $i\in\{1,2\}$, the propositional theory with theorem set
$(D_{t_0})_i$ is essentially incomplete: every consistent extension in
the same language with a computably enumerable axiom set is incomplete.
\end{corollary}
\begin{proof}
Let $U\subseteq S_i$ be the deductive closure of such an axiom set.
The finite proof calculus makes $U$ computably enumerable, and
$U$ contains $(D_{t_0})_i$.
Set $T_i=U$ and $T_j=(D_{t_0})_j$ for $j\ne i$.
This family is consistent. Since $D_{t_0}$ is a subfunctor and
$T_i$ contains its $i$th component, the family is compatible.
Theorem~\ref{thm:cosheaf-observation} therefore makes $U$ incomplete.
\end{proof}

\subsection{First-order theories in the functor category}
\label{subsec:fork-first-order}

Let $N_\Delta=\Delta\N$ be the constant functor on $J$ with value
$\N$ and identity transition morphisms. We use $N_\Delta$ as the
object of variables. For an object $D$ of $\mathcal E$, consider
the first-order language with equality whose only nonlogical symbols
are the propositional atoms specified by $D$. Its terms are variables.

\begin{proposition}\label{prop:fork-first-order-realization}
For every $D\in\mathcal E$, this language has the syntax,
substitution, and initial derivation algebras of
Appendix~\ref{app:syntax-theories}. The syntax and derivation
constructions are stable under parameters in $\mathcal E$.
\end{proposition}
\begin{proof}
The formula constructor endofunctor is
\[
 K_D(Y)=N_\Delta^2+D+Y+Y^2+N_\Delta\times Y.
\]
For $f:Y\to Y'$, set
$K_D(f)=1_{N_\Delta^2}+1_D+f+(f\times f)+(1_{N_\Delta}\times f)$.
Its five summands represent equality, atoms, negation, implication,
and universal quantification. At each component, finite numerical
tree codes form a c.e.\ carrier. Membership in the label carrier
$D_j$ is semidecidable.
For $i=1,2$, relabel each atomic label by $D(0\to i)$ and retain
the variable names and constructors. This defines the transition
morphisms of a formula object $F_D$.

Given an algebra $\alpha:K_D(Y)\to Y$, evaluate its finite trees
from the leaves to the root using the components of $\alpha$.
This is partial recursive and total on each formula carrier.
Naturality of $\alpha$ makes evaluation commute with transitions,
and induction on the trees gives the algebra equation and
uniqueness. If an algebra has a parameter in an object $Z$ of
$\mathcal E$, retain that parameter at every evaluation step.
The same induction proves the parameterized initiality equation
and uniqueness in $\mathcal E/Z$.

Finite lists of variables and their operations are formed
componentwise, with identity transitions on variable names.
Equality of variables is decidable because the variable object
is $N_\Delta$. The occurrence-list and substitution constructions
of Proposition~\ref{prop:categorical-substitution} therefore apply.
Their recursive clauses commute with transitions because relabelling
propositional atoms preserves their empty occurrence lists and the
transition morphisms fix variable names. In particular, empty occurrence lists
define the sentence subobject of $F_D$.

Write $S_D\hookrightarrow F_D$ for the sentence inclusion and let
$A\hookrightarrow S_D$ be an axiom subobject. Use the seven schemes
and two rules of Appendix~\ref{app:first-order-derivations}. At component
$j$, parse a proposed derivation and its rule parameters. Check the
tree form, premise codes, and substitution and free-variable conditions.
For each propositional atom in a formula parameter, test whether its
label $d$ belongs to $D_j$, including at logical axiom leaves. For each
nonlogical axiom leaf also test its sentence for membership in $A_j$.
These finitely many semidecisions succeed together exactly on the
derivation carrier, so that carrier is computably enumerable. Relabelling preserves
each logical scheme and rule and carries nonlogical axiom leaves to nonlogical
axiom leaves because the axioms form a subfunctor.

Evaluation of derivation trees in an algebra for the inference
endofunctor is total on their carriers. Its algebra equation is
proved by induction on the last rule. Naturality of the algebra
operations makes the evaluations a natural transformation.
Induction also gives uniqueness. For $Z\in\mathcal E$ and an axiom
subobject $A\hookrightarrow Z\times S_D$, the conclusion of a tree
at component $j$ is $(z,\varphi)\in Z_j\times (F_D)_j$. Its nodes
retain $z$, and its nonlogical leaves test $(z,\psi)\in A_j$.
Thus the conclusion morphism $c_{Z,A}:P_{Z,A}\to Z\times F_D$
and the inference algebra lie over $Z\times F_D$. For any other
algebra with carrier $h:Y\to Z\times F_D$, evaluation is the unique
algebra morphism $u:P_{Z,A}\to Y$ with $hu=c_{Z,A}$. This gives the initial algebra for the inference endofunctor on
$\mathcal E/(Z\times F_D)$, equivalently on the slice of $\mathcal E/Z$
over $\pi_1:Z\times F_D\to Z$.
\end{proof}

Take $D=V$. Write $F_{\mathrm f}$ for the formula object,
$s_{\mathrm f}:S_{\mathrm f}\hookrightarrow F_{\mathrm f}$ for
its sentence inclusion, and $\theta_{\mathrm f}:V\to S_{\mathrm f}$
for its atomic constructor. The propositional formula
object is $S$. Preserving atoms, negation, and implication defines
a monomorphism $j:S\hookrightarrow S_{\mathrm f}$.
Let $T_{\mathrm p}=D_{t_0}\hookrightarrow S$ be the
propositional theory of Theorem~\ref{thm:fork-proof-presentation}.
Let $T_{\mathrm f}^F\hookrightarrow F_{\mathrm f}$ be the
formula theory generated by the positive atoms from $A$ and the
negative atoms from $B$, and put
$T_{\mathrm f}=s_{\mathrm f}^*T_{\mathrm f}^F$.

\begin{proposition}\label{prop:fork-first-order-conservative}
The theory $T_{\mathrm f}$ is consistent and essentially
incomplete in $\mathcal A$. Its restriction to the propositional
language is $T_{\mathrm p}$, that is,
\[
 j^*T_{\mathrm f}=T_{\mathrm p}.
\]
At either endpoint $i=1,2$, every consistent extension of
$(T_{\mathrm f})_i$ in its first-order language with a c.e.\ axiom set
is incomplete.
\end{proposition}
\begin{proof}
The components
$(0,0,0,0)$, $(0,0)$, and $(0,0)$ give a global element of $V$.
Let $p_*$ be its propositional atom and put
$\top=(p_*\Rightarrow p_*)$. Define a morphism
$r:F_{\mathrm f}\to S$ by the following clauses, writing $P_v$
for the first-order atom with label $v$:
\[
\begin{aligned}
 r(P_v)&=\theta(v),& r(x=y)&=\top,\\
 r(\neg\varphi)&=\nu r(\varphi),&
 r(\varphi\Rightarrow\psi)&=r(\varphi)\Rightarrow r(\psi),\\
 r(\forall x\,\varphi)&=r(\varphi).
\end{aligned}
\]
The clauses commute with transitions and satisfy
$r s_{\mathrm f}j=1_S$. Induction on formulas gives
\[
 r(\varphi[x:=y])=r(\varphi)
\]
whenever the substitution is admissible.

The subobject $r^*T_{\mathrm p}$ contains the generating positive
and negative atoms. It is closed under modus ponens, and erasing
quantification makes it closed under generalization. The first
three logical schemes are sent to the corresponding
propositional schemes. The remaining four schemes are sent to
\[
 H\Rightarrow H,\qquad
 (H\Rightarrow K)\Rightarrow(H\Rightarrow K),\qquad
 \top,\qquad
 \top\Rightarrow(H\Rightarrow H).
\]
Each has a finite propositional derivation, uniformly in $H,K$.
For example, the first two propositional schemes give
\[
\begin{aligned}
 &H\Rightarrow((H\Rightarrow H)\Rightarrow H),\\
 &H\Rightarrow(H\Rightarrow H),\\
 &\bigl(H\Rightarrow((H\Rightarrow H)\Rightarrow H)\bigr)
   \Rightarrow
   \bigl((H\Rightarrow(H\Rightarrow H))
                   \Rightarrow(H\Rightarrow H)\bigr).
\end{aligned}
\]
Two applications of modus ponens yield $H\Rightarrow H$.
Substituting $H\Rightarrow K$ and $p_*$ for $H$ gives the second
formula and $\top$, respectively. The first
propositional scheme and modus ponens applied to $H\Rightarrow H$ give
$\top\Rightarrow(H\Rightarrow H)$.
Leastness of the generated formula theory proves
\[
 T_{\mathrm f}^F\le r^*T_{\mathrm p}.
\]
The inclusion $j$ preserves the signed axioms, the three propositional
schemes, and modus ponens, so $T_{\mathrm p}\le j^*T_{\mathrm f}$.
Put $k=r s_{\mathrm f}:S_{\mathrm f}\to S$. Restricting the formula
inclusion to sentences gives $T_{\mathrm f}\le k^*T_{\mathrm p}$.
Since $kj=1_S$, this gives $j^*T_{\mathrm f}\le T_{\mathrm p}$.
Writing $\nu_{\mathrm f}$ for sentence negation, we have
$k\nu_{\mathrm f}=\nu k$, and hence
\[
 T_{\mathrm f}\wedge\nu_{\mathrm f}^*T_{\mathrm f}
 \le k^*(T_{\mathrm p}\wedge\nu^*T_{\mathrm p})=\bot .
\]

Let $U\hookrightarrow S_{\mathrm f}$ be a consistent closed extension of
$T_{\mathrm f}$. The positive and negative axioms give
\[
 A\leq \theta_{\mathrm f}^{*}U,\qquad
 B\leq \theta_{\mathrm f}^{*}\nu_{\mathrm f}^{*}U .
\]
If $U$ were complete, $U$ and $\nu_{\mathrm f}^{*}U$ would be
complementary. Their inverse images under $\theta_{\mathrm f}$
would separate $A$ and $B$, contrary to their inseparability.
This is Corollary~\ref{cor:domain-incompleteness} in $\mathcal A$.

For the endpoint assertion, let $U_i$ be the deductive closure
of a consistent c.e.\ set of first-order axioms, and suppose
$(T_{\mathrm f})_i\subseteq U_i$. The inverse image $j_i^*U_i$ is a
consistent c.e.\ propositional theory extending
$(T_{\mathrm p})_i$. Completeness of $U_i$ would imply
completeness of $j_i^*U_i$, since $j_i$ preserves negation.
This contradicts Corollary~\ref{cor:cosheaf-essential-incompleteness}.
\end{proof}

\subsection{Categories associated with theory extensions}

\begin{corollary}\label{cor:fork-associated-theories}
Let $t:T_{\mathrm f}\hookrightarrow S_{\mathrm f}$ be the
first-order theory above, and let
$u:U\hookrightarrow S_{\mathrm f}$ be any consistent internal
sentence theory with $t=ui$ for a monomorphism
$i:T_{\mathrm f}\hookrightarrow U$. Then
\[
 \mathcal A[T_{\mathrm f}]=\Par(\mathcal E/T_{\mathrm f},\Mono),
 \qquad
 \mathcal A[U]=\Par(\mathcal E/U,\Mono)
\]
are Turing categories, and postcomposition gives the fully
faithful functor
\[
 \Sigma_i:\mathcal A[T_{\mathrm f}]\longrightarrow\mathcal A[U]
\]
of Theorem~\ref{thm:theory-extension}. It is not essentially
surjective when $i$ is not an isomorphism. The theory
\[
 1_U\times u:U\times U\hookrightarrow U\times S_{\mathrm f}
\]
is essentially incomplete in $\mathcal A[U]$. Its theorem and
refutation domains are inseparable there.
\end{corollary}
\begin{proof}
Propositions~\ref{prop:fork-total} and
\ref{prop:fork-structure} identify $\mathcal E$ as regular and
extensive and $\mathcal A$ as its split partial-map category.
Proposition~\ref{prop:fork-recursion} identifies
$\mathcal A\simeq\Split(\mathscr R_X)$.
Proposition~\ref{prop:fork-first-order-realization} supplies the
parameter-stable syntax and deduction required by the general
theorem. The global sentence $j(\top)$ is a logical theorem, so
there is a total morphism $\lambda_U:1\to U$ with $u\lambda_U=j(\top)$.
Proposition~\ref{prop:fork-first-order-conservative} and
Corollary~\ref{cor:domain-incompleteness} show that every such
$u$ is essentially incomplete and that its theorem and
refutation domains are inseparable in $\mathcal A$. If the constant
domains had a separator in the slice, its inverse image along
$\langle \lambda_U !_{S_{\mathrm f}},1_{S_{\mathrm f}}\rangle:
 S_{\mathrm f}\to U\times S_{\mathrm f}$ would separate $u$ and
$\nu_{\mathrm f}^*u$ in $\mathcal A$.
Theorems~\ref{thm:theory-extension} and
\ref{thm:associated-incompleteness} give the assertions about the
comparison functor and essential incompleteness.
\end{proof}
